\section{LangChain 实现 ReAct Agent}

\subsection{Prompt 模板设计}

LangChain 中 \texttt{create\_react\_agent} 使用了一个标准化的 Prompt 模板，其设计强调\textbf{泛化性和通用性}：

\begin{quote}
\textit{Answer the following questions as best you can. You have access to the following tools: \{tools\}. Use the following format: Question / Thought / Action / Action Input / Observation / ... / Final Answer. Begin! Question: \{input\} \{agent\_scratchpad\}}
\end{quote}

模板包含两个动态占位符：
\begin{itemize}
    \item \texttt{tools}：可用的工具列表
    \item \texttt{agent\_scratchpad}：交互过程中动态维护的历史记录
\end{itemize}

\subsection{三步骤构建流程}

\begin{enumerate}
    \item 声明语言模型（如 GPT-4.1 Nano）
    \item 调用 \texttt{create\_react\_agent} 创建 Agent
    \item 将 Agent 放入 \texttt{AgentExecutor} 中执行
\end{enumerate}

\subsection{工作机制：基于文本模板的迭代}

LangChain 的 ReAct 实现有一个重要特征——它\textbf{不使用多轮消息（message list）}，而是将所有历史信息不断追加到\textbf{同一个 human 消息}中：

\begin{warningbox}{LangChain ReAct 的历史管理方式}
每一轮新的交互都是一次新的 LLM 调用。历史信息（Thought、Action、Observation）被追加到 \texttt{agent\_scratchpad} 中，作为单条 human 消息的一部分发送给模型。这意味着没有多轮对话的 message 结构，所有上下文都挤在一个消息里。
\end{warningbox}

\subsection{本章小结}
LangChain 的 ReAct 实现基于 Prompt 模板 + 文本追加的方式管理交互历史，结构简单但缺乏多轮消息的清晰组织。

